\section{Limitations and Discussion}
Despite the progress of MemT-1.0 in native memory, it still has several limitations.
The fixed-size memory state imposes a capacity constraint.
For the fixed memory, MemT-1.0 maintains an $O(1)$ state with respect to the interaction history length $N$.
Compared with methods whose stored representations grow as $O(N)$, this bounded state has limited capacity to preserve detailed historical information.
As interactions accumulate, new information must be incorporated into an existing representation, creating trade-offs between retaining earlier details and adding recent information.
The resulting loss of memory fidelity remains a challenge for tasks requiring precise recall.
At the same time, a fixed-size state offers opportunities beyond reducing storage.
Each update incorporates new information into a representation, which has already been shaped by earlier interactions.
This enables successive compression of accumulated history.
Compared with maintaining a growing collection of historical records, this recurrent organization may provide a basis for latent reflection.
It potentially allows iterative consolidation of experience within the memory state.
Such a mechanism could support online self-evolution by allowing accumulated experience to influence subsequent behavior.

Finally, the relationship between local memory and hyper memory also deserves further study.
Local memory represents information accumulated during interaction, while hyper memory encodes learned capabilities for storing and utilizing that information.
An important question is whether experience accumulated in local memory can be consolidated into hyper-memory parameters through further training, and how knowledge encoded in hyper memory can be instantiated and adapted through local memory.
Therefore, memory capability itself can be viewed as a higher-order form of learned information.
Exploring this relationship may help connect adaptation within individual interactions with improvements in memory capability across tasks.